% ============================================================
%  ARS TYPOGRAPHICA — Template LaTeX
%  Charte Graphique Minimaliste, inspiration TAOCP / Knuth
%  Version 1.1 — Mars 2026
%  Inclus : environnements mathématiques complets numérotés
% ============================================================
%
%  Compilez avec : pdflatex (2×) ou lualatex (recommandé)
%
%  Dépendances :
%    geometry, fontenc, inputenc, microtype,
%    booktabs, xcolor, listings, hyperref,
%    titlesec, fancyhdr, parskip,
%    amsmath, amssymb, amsthm, thmtools,
%    mdframed, enumitem, tcolorbox
% ============================================================

\documentclass[11pt, a4paper, twoside]{article}

% ─── ENCODAGE & LANGUE ───────────────────────────────────────
\usepackage[T1]{fontenc}
\usepackage[utf8]{inputenc}
\usepackage[french]{babel}

% ─── TYPOGRAPHIE ─────────────────────────────────────────────
\usepackage{lmodern}         % Latin Modern (amélioré de Computer Modern)
\usepackage{microtype}       % Crénage et expansion optimaux
% \usepackage{inconsolata}   % Décommenter pour Inconsolata (monospace)

% ─── MISE EN PAGE ────────────────────────────────────────────
\usepackage[
  a4paper,
  top    = 30mm,
  bottom = 28mm,
  left   = 28mm,
  right  = 28mm,
  headsep= 8mm,
  footskip = 12mm,
]{geometry}

% ─── COULEURS ────────────────────────────────────────────────
\usepackage[dvipsnames]{xcolor}

\definecolor{ink}{HTML}{111111}
\definecolor{darkgray}{HTML}{333333}
\definecolor{midgray}{HTML}{666666}
\definecolor{lightgray}{HTML}{AAAAAA}
\definecolor{rulegray}{HTML}{CCCCCC}
\definecolor{codebg}{HTML}{F5F5F5}
\definecolor{thmbg}{HTML}{F5F5F5}

\color{ink}

% ─── MATHÉMATIQUES ───────────────────────────────────────────
\usepackage{amsmath}
\usepackage{amssymb}
\usepackage{amsthm}
\usepackage{mathtools}

% ─── CADRES POUR THÉORÈMES ───────────────────────────────────
\usepackage[framemethod=TikZ]{mdframed}

% Style cadre léger (définitions, théorèmes)
\mdfdefinestyle{thmstyle}{
  linecolor      = rulegray,
  linewidth      = 0.4pt,
  backgroundcolor= thmbg,
  innertopmargin = 8pt,
  innerbottommargin = 8pt,
  innerleftmargin = 10pt,
  innerrightmargin = 10pt,
  skipabove      = 10pt,
  skipbelow      = 8pt,
  nobreak        = true,
}

% Style cadre trait gauche (exemples, remarques)
\mdfdefinestyle{examplestyle}{
  topline        = false,
  bottomline     = false,
  rightline      = false,
  linecolor      = rulegray,
  linewidth      = 1.2pt,
  backgroundcolor= white,
  innertopmargin = 4pt,
  innerbottommargin = 4pt,
  innerleftmargin = 10pt,
  innerrightmargin = 0pt,
  skipabove      = 8pt,
  skipbelow      = 6pt,
}

% Style encadré exercice
\mdfdefinestyle{exercisestyle}{
  linecolor      = ink,
  linewidth      = 0.4pt,
  backgroundcolor= white,
  innertopmargin = 8pt,
  innerbottommargin = 8pt,
  innerleftmargin = 10pt,
  innerrightmargin = 10pt,
  skipabove      = 10pt,
  skipbelow      = 8pt,
}

% ─── COMPTEURS ET STYLES DE THÉORÈMES ────────────────────────
%
%  Convention de numérotation :
%  - Définition, Théorème, Proposition, Lemme, Corollaire
%    → compteur unique par section : Théorème 1.1, Théorème 2.3 …
%  - Exemple, Remarque → compteur propre par section
%  - Exercice → compteur global par section, numéroté § 1, § 2 …
%  - Correction → même numéro que l'exercice correspondant
%  - Preuve → non numérotée, terminée par ■

\newcounter{thmcnt}[section]
\renewcommand{\thethmcnt}{\thesection.\arabic{thmcnt}}

\newcounter{examplecnt}[section]
\renewcommand{\theexamplecnt}{\thesection.\arabic{examplecnt}}

\newcounter{exercisecnt}[section]
\renewcommand{\theexercisecnt}{\thesection.\arabic{exercisecnt}}

% Macro interne pour les environnements encadrés
\newcommand{\thmheader}[2]{%
  \textbf{#1\,\thethmcnt}%
  \ifx\relax#2\relax\else~---~\textit{#2}\fi
  \stepcounter{thmcnt}%
}

% ── DÉFINITION ───────────────────────────────────────────────
\newenvironment{definition}[1][\relax]{%
  \refstepcounter{thmcnt}%
  \begin{mdframed}[style=thmstyle]
  \noindent\textbf{Définition\,\thethmcnt}%
  \ifx\relax#1\relax\else~\textup{(#1)}\fi.\quad\ignorespaces
}{%
  \end{mdframed}
}

% ── THÉORÈME ─────────────────────────────────────────────────
\newenvironment{theorem}[1][\relax]{%
  \refstepcounter{thmcnt}%
  \begin{mdframed}[style=thmstyle]
  \noindent\textbf{Théorème\,\thethmcnt}%
  \ifx\relax#1\relax\else~\textup{(#1)}\fi.\quad\ignorespaces
}{%
  \end{mdframed}
}

% ── PROPOSITION ──────────────────────────────────────────────
\newenvironment{proposition}[1][\relax]{%
  \refstepcounter{thmcnt}%
  \begin{mdframed}[style=thmstyle]
  \noindent\textbf{Proposition\,\thethmcnt}%
  \ifx\relax#1\relax\else~\textup{(#1)}\fi.\quad\ignorespaces
}{%
  \end{mdframed}
}

% ── LEMME ────────────────────────────────────────────────────
\newenvironment{lemma}[1][\relax]{%
  \refstepcounter{thmcnt}%
  \begin{mdframed}[style=thmstyle]
  \noindent\textbf{Lemme\,\thethmcnt}%
  \ifx\relax#1\relax\else~\textup{(#1)}\fi.\quad\ignorespaces
}{%
  \end{mdframed}
}

% ── COROLLAIRE ───────────────────────────────────────────────
\newenvironment{corollary}[1][\relax]{%
  \refstepcounter{thmcnt}%
  \begin{mdframed}[style=thmstyle]
  \noindent\textbf{Corollaire\,\thethmcnt}%
  \ifx\relax#1\relax\else~\textup{(#1)}\fi.\quad\ignorespaces
}{%
  \end{mdframed}
}

% ── PREUVE ───────────────────────────────────────────────────
% Non numérotée. Se termine par le carré ■ (QED)
\renewenvironment{proof}[1][Démonstration]{%
  \noindent\textit{#1.}\quad\ignorespaces
}{%
  \hfill$\blacksquare$\par\vspace{4pt}
}

% ── EXEMPLE ──────────────────────────────────────────────────
\newenvironment{example}[1][\relax]{%
  \refstepcounter{examplecnt}%
  \begin{mdframed}[style=examplestyle]
  \noindent\textbf{Exemple\,\theexamplecnt}%
  \ifx\relax#1\relax\else~\textup{(#1)}\fi.\quad\ignorespaces
}{%
  \end{mdframed}
}

% ── REMARQUE ─────────────────────────────────────────────────
\newenvironment{remark}[1][\relax]{%
  \refstepcounter{examplecnt}%
  \begin{mdframed}[style=examplestyle]
  \noindent\textbf{Remarque\,\theexamplecnt}%
  \ifx\relax#1\relax\else~\textup{(#1)}\fi.\quad\ignorespaces\itshape
}{%
  \end{mdframed}
}

% ── EXERCICE ─────────────────────────────────────────────────
\newenvironment{exercise}[1][\relax]{%
  \refstepcounter{exercisecnt}%
  \begin{mdframed}[style=exercisestyle]
  \noindent\textbf{Exercice\,\theexercisecnt}%
  \ifx\relax#1\relax\else~\textup{---~#1}\fi.\quad\ignorespaces
}{%
  \end{mdframed}
}

% ── CORRECTION ───────────────────────────────────────────────
% Même compteur que l'exercice (ne l'incrémente pas)
\newenvironment{correction}[1][\relax]{%
  \begin{mdframed}[style=examplestyle]
  \noindent\textit{Correction de l'exercice~\theexercisecnt}%
  \ifx\relax#1\relax\else~\textup{(#1)}\fi.\quad\ignorespaces
}{%
  \hfill$\lozenge$%
  \end{mdframed}
}

% ─── RÉFÉRENCES CROISÉES ─────────────────────────────────────
% Utilisation : \label{thm:moi} dans l'env, puis \refthm{thm:moi}
\newcommand{\refthm}[1]{Théorème~\ref{#1}}
\newcommand{\refprop}[1]{Proposition~\ref{#1}}
\newcommand{\reflem}[1]{Lemme~\ref{#1}}
\newcommand{\refcor}[1]{Corollaire~\ref{#1}}
\newcommand{\refdef}[1]{Définition~\ref{#1}}
\newcommand{\refex}[1]{Exemple~\ref{#1}}
\newcommand{\refexo}[1]{Exercice~\ref{#1}}

% ─── NOTATIONS MATHÉMATIQUES ─────────────────────────────────
\newcommand{\N}{\mathbb{N}}
\newcommand{\Z}{\mathbb{Z}}
\newcommand{\Q}{\mathbb{Q}}
\newcommand{\R}{\mathbb{R}}
\newcommand{\C}{\mathbb{C}}
\newcommand{\K}{\mathbb{K}}
\newcommand{\E}{\mathbb{E}}
\newcommand{\Prob}{\mathbb{P}}
\newcommand{\Var}{\operatorname{Var}}
\newcommand{\Cov}{\operatorname{Cov}}
\newcommand{\1}{\mathbf{1}}
\newcommand{\norm}[1]{\left\|#1\right\|}
\newcommand{\abs}[1]{\left|#1\right|}
\newcommand{\floor}[1]{\left\lfloor#1\right\rfloor}
\newcommand{\ceil}[1]{\left\lceil#1\right\rceil}
\DeclareMathOperator{\Card}{Card}
\DeclareMathOperator{\Im}{Im}
\DeclareMathOperator{\Ker}{Ker}
\DeclareMathOperator{\rk}{rk}
\DeclareMathOperator{\tr}{tr}
\DeclareMathOperator*{\argmin}{arg\,min}
\DeclareMathOperator*{\argmax}{arg\,max}

% ─── EN-TÊTES ET PIEDS ───────────────────────────────────────
\usepackage{fancyhdr}
\pagestyle{fancy}
\fancyhf{}
\fancyhead[LE]{\small\itshape\color{lightgray}\leftmark}
\fancyhead[RO]{\small\itshape\color{lightgray}\rightmark}
\renewcommand{\headrulewidth}{0.4pt}
\renewcommand{\headrule}{\color{rulegray}\hrule width\headwidth height\headrulewidth}
\renewcommand{\footrulewidth}{0pt}
\fancyfoot[C]{\small\color{midgray}\thepage}
\fancypagestyle{plain}{\fancyhf{}\fancyfoot[C]{\small\color{midgray}\thepage}\renewcommand{\headrulewidth}{0pt}}

% ─── TITRES DE SECTIONS ──────────────────────────────────────
\usepackage{titlesec}
\titleformat{\section}{\fontsize{17}{22}\bfseries\color{ink}}{\S\,\thesection}{1em}{}
\titlespacing*{\section}{0pt}{24pt}{12pt}
\titleformat{\subsection}{\fontsize{13}{18}\bfseries\color{ink}}{\S\,\thesubsection}{1em}{}
\titlespacing*{\subsection}{0pt}{18pt}{8pt}
\titleformat{\subsubsection}{\fontsize{11}{16}\bfseries\itshape\color{ink}}{\S\,\thesubsubsection}{1em}{}
\titlespacing*{\subsubsection}{0pt}{12pt}{6pt}

% ─── PARAGRAPHES ─────────────────────────────────────────────
\usepackage{parskip}
\setlength{\parskip}{6pt}
\setlength{\parindent}{0pt}

% ─── TABLEAUX ────────────────────────────────────────────────
\usepackage{booktabs}
\usepackage{array}
\usepackage{longtable}

\newenvironment{chartable}[3]{%
  \begin{table}[htbp]\caption{#2}\label{#3}\centering\small\begin{tabular}{#1}\toprule
}{%
  \bottomrule\end{tabular}\end{table}
}

% ─── CODE SOURCE ─────────────────────────────────────────────
\usepackage{listings}
\lstset{
  basicstyle=\small\ttfamily\color{ink},
  backgroundcolor=\color{codebg},
  frame=single, framesep=4pt,
  rulecolor=\color{rulegray}, framerule=0.3pt,
  numbers=left, numberstyle=\tiny\color{lightgray}, numbersep=8pt,
  breaklines=true, showstringspaces=false, tabsize=4,
  keywordstyle=\bfseries,
  commentstyle=\itshape\color{midgray},
  stringstyle=\ttfamily,
  literate={<-}{$\leftarrow$}2 {->}{$\rightarrow$}2
           {<=}{$\leq$}2 {>=}{$\geq$}2 {!=}{$\neq$}2,
}
\newcommand{\code}[1]{\texttt{#1}}

% ─── FIGURES ─────────────────────────────────────────────────
\usepackage{graphicx}
\usepackage{caption}
\captionsetup{font={small,it}, labelfont={small,bf}, labelsep=period, textcolor=midgray}

% ─── LIENS ───────────────────────────────────────────────────
\usepackage[colorlinks=true, linkcolor=ink, urlcolor=darkgray, citecolor=darkgray,
            pdftitle={Document — Ars Typographica}]{hyperref}

% ─── MACROS UTILITAIRES ──────────────────────────────────────
\newcommand{\ksec}[1]{\S\,#1}
\newcommand{\lightrule}{{\color{rulegray}\rule{\linewidth}{0.4pt}}\vspace{4pt}}
\newcommand{\heavyrule}{{\color{ink}\rule{\linewidth}{1.2pt}}\vspace{6pt}}

\renewenvironment{abstract}{%
  \small\itshape\noindent\textbf{Résumé.} ---
}{\par\vspace{6pt}\lightrule\vspace{6pt}}

% ─── PAGE DE TITRE ───────────────────────────────────────────
\renewcommand{\maketitle}{%
  \begin{center}
    \heavyrule
    {\fontsize{22}{28}\bfseries\@title\par}
    \vspace{4pt}
    {\large\itshape\@author\par}
    \vspace{2pt}
    {\small\color{midgray}\@date\par}
    \vspace{6pt}
    \lightrule
  \end{center}
  \vspace{8pt}
}

% ============================================================
%  DÉBUT DU DOCUMENT EXEMPLE
% ============================================================
\title{Titre du Document}
\author{Prénom Nom}
\date{Mars 2026}

\begin{document}

\maketitle

\begin{abstract}
Résumé concis du document en italique. Ce texte donne les grandes lignes
du contenu et des conclusions, en deux à cinq phrases.
\end{abstract}

\tableofcontents
\vspace{12pt}\lightrule

% ============================================================
%  § 1 — DÉFINITIONS ET RÉSULTATS PRINCIPAUX
% ============================================================
\section{Définitions et résultats principaux}

Le texte courant est composé en Computer Modern 11\,pt, justifié,
interlignage 16\,pt. Les variables sont en italique : $f\colon\R^n\to\R$,
la norme euclidienne est $\norm{x} = \sqrt{x^\top x}$.

\begin{definition}[Convergence uniforme]\label{def:cvunif}
  Une suite de fonctions $(f_n)$ converge \emph{uniformément} vers $f$
  sur $A \subset \R$ si
  \[
    \forall \varepsilon > 0,\;\exists N\in\N,\;\forall n\geq N,\;
    \forall x\in A,\quad \abs{f_n(x)-f(x)} < \varepsilon.
  \]
\end{definition}

\begin{theorem}[Borne inférieure du tri]\label{thm:triling}
  Tout algorithme de tri par comparaison requiert $\Omega(n\log n)$
  comparaisons dans le pire cas.
\end{theorem}

\begin{proof}
  L'arbre de décision sur $n$ éléments possède au moins $n!$ feuilles,
  donc sa hauteur est au moins $\log_2(n!) = \Omega(n\log n)$.
\end{proof}

\begin{proposition}[Critère de Cauchy]\label{prop:cauchy}
  Une suite réelle $(u_n)$ est convergente si et seulement si elle est de
  Cauchy, c'est-à-dire si
  \[
    \forall \varepsilon > 0,\;\exists N\in\N,\;\forall p,q\geq N,\quad
    \abs{u_p - u_q} < \varepsilon.
  \]
\end{proposition}

\begin{lemma}[Sous-suite]\label{lem:sousuite}
  Toute suite convergente est de Cauchy. La réciproque est vraie dans
  $\R$ mais fausse dans $\Q$.
\end{lemma}

\begin{proof}
  Si $u_n \to \ell$ alors, pour $\varepsilon > 0$, il existe $N$ tel que
  $\abs{u_n - \ell} < \varepsilon/2$ pour $n\geq N$; par inégalité
  triangulaire $\abs{u_p - u_q} < \varepsilon$ pour $p,q\geq N$.
\end{proof}

\begin{corollary}\label{cor:cvabsolue}
  Toute série absolument convergente est convergente dans $\R$.
\end{corollary}

% ============================================================
%  § 2 — EXEMPLES ET REMARQUES
% ============================================================
\section{Exemples et remarques}

\begin{example}[Suite géométrique]\label{ex:geom}
  La suite $u_n = r^n$ converge vers $0$ si $\abs{r}<1$, vers $1$ si
  $r=1$, et diverge sinon. Le cas $r = 0.9$ donne
  \[
    u_{100} = 0.9^{100} \approx 2.656 \times 10^{-5}.
  \]
\end{example}

\begin{example}[KMP en pratique]\label{ex:kmp}
  L'algorithme de Knuth–Morris–Pratt recherche un motif de longueur $m$
  dans un texte de longueur $n$ en $O(n+m)$ en prétraitant le tableau
  \emph{lps} (longest proper prefix-suffix) du motif en $O(m)$.
\end{example}

\begin{remark}
  La complexité $O(n+m)$ de KMP est optimale au sens de la complexité en
  comparaisons de caractères, à une constante multiplicative près.
\end{remark}

\begin{remark}[Sur la \refdef{def:cvunif}]
  La convergence uniforme implique la convergence simple, mais la
  réciproque est fausse. La suite $f_n(x) = x^n$ sur $[0,1]$ converge
  simplement vers $\mathbf{1}_{\{1\}}$, mais non uniformément.
\end{remark}

% ============================================================
%  § 3 — EXERCICES ET CORRECTIONS
% ============================================================
\section{Exercices}

\begin{exercise}[Convergence uniforme vs simple]
  Montrer que la suite $f_n(x) = \dfrac{x}{1+nx^2}$ converge
  uniformément vers $0$ sur $\R$.
\end{exercise}

\begin{correction}
  On a $f_n(x) = \dfrac{x}{1+nx^2}$. Par l'inégalité AM-GM,
  $1+nx^2 \geq 2\sqrt{n}\,\abs{x}$, donc
  \[
    \abs{f_n(x)} \leq \frac{\abs{x}}{2\sqrt{n}\,\abs{x}} = \frac{1}{2\sqrt{n}}
    \xrightarrow[n\to\infty]{} 0
  \]
  uniformément en $x \in \R^*$; la borne est valable aussi en $x=0$.
\end{correction}

\begin{exercise}[Complexité du tri par fusion]
  Établir la récurrence $T(n) = 2\,T(n/2) + \Theta(n)$ du tri par fusion
  et en déduire $T(n) = \Theta(n\log n)$ par le théorème maître.
\end{exercise}

\begin{correction}
  L'appel récursif sur deux moitiés donne le facteur $2\,T(n/2)$.
  La fusion de deux listes triées de taille $n/2$ coûte $\Theta(n)$.
  Le théorème maître (cas 2, $f(n) = \Theta(n^{\log_2 2}) = \Theta(n)$)
  donne directement $T(n) = \Theta(n\log n)$.
\end{correction}

\begin{exercise}[Inégalité de Jensen]
  Soit $\varphi\colon\R\to\R$ convexe et $X$ une variable aléatoire
  intégrable. Démontrer que $\varphi(\E[X]) \leq \E[\varphi(X)]$.
\end{exercise}

% ============================================================
%  § 4 — CODE ET TABLEAUX
% ============================================================
\section{Code et tableaux}

\begin{lstlisting}[language=Python, caption={Tableau lps de KMP}]
def compute_lps(pattern):
    """Compute longest proper prefix-suffix array."""
    m, lps = len(pattern), [0] * len(pattern)
    length = 0; i = 1
    while i < m:
        if pattern[i] == pattern[length]:
            length += 1; lps[i] = length; i += 1
        elif length != 0:
            length = lps[length - 1]
        else:
            lps[i] = 0; i += 1
    return lps
\end{lstlisting}

\begin{chartable}{lrrrr}{Complexités — algorithmes de tri}{tab:tri}
  \textbf{Algorithme} & \textbf{Meilleur} & \textbf{Moyen}
    & \textbf{Pire} & \textbf{Mémoire} \\
  \midrule
  Insertion sort & $O(n)$ & $O(n^2)$ & $O(n^2)$ & $O(1)$ \\
  Mergesort & $O(n\log n)$ & $O(n\log n)$ & $O(n\log n)$ & $O(n)$ \\
  Quicksort & $O(n\log n)$ & $O(n\log n)$ & $O(n^2)$ & $O(\log n)$ \\
  Heapsort & $O(n\log n)$ & $O(n\log n)$ & $O(n\log n)$ & $O(1)$ \\
  Timsort & $O(n)$ & $O(n\log n)$ & $O(n\log n)$ & $O(n)$ \\
\end{chartable}

% ─── RÉFÉRENCES ──────────────────────────────────────────────
\begin{thebibliography}{9}
  \bibitem{knuth1997}
    D.\,E.~Knuth, \textit{The Art of Computer Programming, Vol.\,1},
    3\textsuperscript{e}~éd., Addison-Wesley, 1997.
  \bibitem{knuth1998}
    D.\,E.~Knuth, \textit{The Art of Computer Programming, Vol.\,3},
    2\textsuperscript{e}~éd., Addison-Wesley, 1998.
  \bibitem{rudin}
    W.~Rudin, \textit{Principles of Mathematical Analysis},
    3\textsuperscript{e}~éd., McGraw-Hill, 1976.
\end{thebibliography}

\end{document}
% ============================================================
%  FIN — ARS TYPOGRAPHICA v1.1
% ============================================================
